%-------------------------
% Resume in Latex
% Based off of: https://github.com/sb2nov/resume
% License : MIT
%
% This file is the rendering shell for the Resume Builder. Everything in the
% preamble (packages, margins, \resume* custom commands) is fixed; only the
% body sections between \begin{document} and \end{document} are produced from
% structured data by backend/resume_render.py. Do NOT hand-edit the {{...}}
% placeholders — they are filled in programmatically.
%------------------------

\documentclass[letterpaper,10pt]{article}

\usepackage{latexsym}
\usepackage[empty]{fullpage}
\usepackage{titlesec}
\usepackage{marvosym}
\usepackage[usenames,dvipsnames]{color}
\usepackage{verbatim}
\usepackage{enumitem}
\usepackage[hidelinks]{hyperref}
\usepackage{fancyhdr}
\usepackage[english]{babel}
\usepackage{tabularx}
% NOTE: fontawesome5 was in the source template but no \fa… icons are used.
% It requires loadable OTF fonts that Tectonic's XeTeX can't resolve in this
% environment (aborts with SIGABRT), so it's intentionally omitted.
\usepackage{multicol}
\usepackage{setspace}
\setstretch{1.02}
\setlength{\multicolsep}{-3.0pt}
\setlength{\columnsep}{-1pt}
% glyphtounicode / \pdfgentounicode are pdfTeX-only primitives used to make the
% PDF ATS-parsable. Tectonic compiles with XeTeX (which emits unicode-mapped
% PDFs natively and does NOT define these), so we load them only when running
% under pdfTeX. Without this guard XeTeX aborts on an undefined control sequence.
\ifdefined\pdfgentounicode
  \input{glyphtounicode}
\fi

\pagestyle{fancy}
\fancyhf{}
\fancyfoot{}
\renewcommand{\headrulewidth}{0pt}
\renewcommand{\footrulewidth}{0pt}

% Optimized margins for professional appearance
\addtolength{\oddsidemargin}{-0.55in}
\addtolength{\evensidemargin}{-0.55in}
\addtolength{\textwidth}{1.1in}
\addtolength{\topmargin}{-0.65in}
\addtolength{\textheight}{1.4in}

\urlstyle{same}

\raggedbottom
\raggedright
\setlength{\tabcolsep}{0in}

% Professional section formatting with consistent spacing.
% The leading \vspace pulls the section title up toward the previous section —
% i.e. it controls the gap BETWEEN sections (Education→Experience, etc.).
% The trailing \vspace (after \titlerule) is the gap below the rule, within the
% section, and is left as-is.
\titleformat{\section}{
  \vspace{-7pt}\scshape\raggedright\large\bfseries
}{}{0em}{}[\color{black}\titlerule \vspace{-4pt}]

% Ensure that generate pdf is machine readable/ATS parsable (pdfTeX only)
\ifdefined\pdfgentounicode \pdfgentounicode=1 \fi

%-------------------------
% Custom commands with balanced spacing
\newcommand{\resumeItem}[1]{
  \item\small{
    {#1 \vspace{-2pt}}
  }
}

\newcommand{\resumeSubheading}[4]{
  \vspace{-2pt}\item
    \begin{tabular*}{1.0\textwidth}[t]{l@{\extracolsep{\fill}}r}
      \textbf{#1} & \textbf{\small #2} \\
      \textit{\small#3} & \textit{\small #4} \\
    \end{tabular*}\vspace{-7pt}
}

\newcommand{\resumeProjectHeading}[2]{
    \item
    \begin{tabular*}{1.001\textwidth}{l@{\extracolsep{\fill}}r}
      \small#1 & \textbf{\small #2}\\
    \end{tabular*}\vspace{-7pt}
}

\renewcommand\labelitemi{$\vcenter{\hbox{\tiny$\bullet$}}$}
\renewcommand\labelitemii{$\vcenter{\hbox{\tiny$\bullet$}}$}

\newcommand{\resumeSubHeadingListStart}{\begin{itemize}[leftmargin=0.0in, label={}]}
\newcommand{\resumeSubHeadingListEnd}{\end{itemize}}
\newcommand{\resumeItemListStart}{\begin{itemize}}
\newcommand{\resumeItemListEnd}{\end{itemize}\vspace{-5pt}}

%-------------------------------------------
%%%%%%  RESUME STARTS HERE  %%%%%%%%%%%%%%%%%%%%%%%%%%%%

\begin{document}

%----------HEADING----------
\begin{center}
    {\Huge \scshape {{FULL_NAME}}} \\ \vspace{3pt}
    \small {{CONTACT_LINE}}
    \vspace{-8pt}
\end{center}

{{EDUCATION_SECTION}}
{{EXPERIENCE_SECTION}}
{{SKILLS_SECTION}}
{{PROJECTS_SECTION}}

\end{document}
